\documentclass[11pt,a4paper]{article}
\usepackage[utf8]{inputenc}
\usepackage[T1]{fontenc}
\usepackage{lmodern}
\usepackage[spanish,es-tabla]{babel}
\usepackage{amsmath,amssymb,mathtools}
\usepackage{booktabs,array,longtable}
\usepackage{graphicx,float}
\usepackage[margin=2.3cm]{geometry}
\usepackage{microtype,caption,xcolor,listings}
\usepackage[hidelinks]{hyperref}
\lstset{basicstyle=\ttfamily\scriptsize,breaklines=true,keywordstyle=\color{blue!60!black},commentstyle=\color{gray},frame=single,numbers=left,numberstyle=\tiny\color{gray}}
\captionsetup{font=small,labelfont=bf}
\title{\textbf{Segundo trabajo de evaluación continua}\\[2pt]\large Renta Fija y sus Derivados -- Estimación de la estructura temporal de los tipos de interés cupón cero}
\author{Andrés Otárola \and Luiz Luzuriaga}
\date{Máster en Ciencias Actuariales y Financieras -- Universidad de Alcalá\\Curso 2026-27 \quad Profesor: Eliseo Navarro}
\begin{document}
\maketitle
\begin{abstract}\noindent A partir de las cotizaciones de bonos y obligaciones del Estado del 11 de septiembre de 2026 (fecha valor 14 de septiembre) se estima la ETTI cupón cero con el modelo de Svensson y, como contraste, con splines cúbicas sobre la función de descuento. Se utilizan 57 títulos (precio ex-cupón medio compra/venta) y se minimizan los errores de precio ponderados por el inverso de la duración. El modelo de Svensson ajusta con un error cuadrático medio de 3{,}9 puntos básicos en TIR y la spline con 2{,}8. La curva es creciente, del 2{,}62\,\% a 0,1 años al 4{,}70\,\% a 45 años.\end{abstract}

\section{Datos y tratamiento}
\paragraph{Muestra.} El fichero de cotizaciones contiene 58 bonos y obligaciones; la Obligación 4,000\,\% 10/2064 no tiene cotización y se excluye, por lo que quedan 57 títulos (2 son cupón cero: el Bono 0,000\,\% 01/2027 y la Obligación 0,000\,\% 01/2028). Se usa el \textbf{precio ex-cupón medio} $P^*=\frac12(P^{compra}+P^{venta})$ en porcentaje del nominal. La fecha de contratación es el 11/09/2026 y la fecha valor el 14/09/2026 (dato del propio fichero, D+1 hábil).
\paragraph{Fechas de pago.} El fichero indica solo mes y año de vencimiento. Se ha tomado el día habitual de las emisiones del Tesoro según el mes (31/01, 30/03, 30/04, 31/05, 30/07, 31/10) y cupones anuales en el aniversario del vencimiento. Como comprobación, con estas fechas, el precio bruto y la convención del Tema~1 la TIR calculada coincide con la TIR publicada (media compra/venta) con una diferencia media absoluta de 0,3 puntos básicos; la única diferencia relevante (8,4\,pb) es la Obligación 1,300\,\% 10/2026, a mes y medio del vencimiento, donde el redondeo del precio se amplifica en el tipo anual.
\paragraph{Plazos y precio bruto.} Según el Tema~1 el cupón corrido es $CC=C\cdot\frac{d}{b}$, con $d$ los días desde el último cupón hasta la fecha valor y $b$ los días entre cupones (ACT/ACT), y el precio bruto es $P=P^*+CC$. Los flujos futuros se sitúan en $t_k=\tau+k$ años, con $\tau=\frac{\text{días hasta el próximo cupón}}{b}$ y $k=0,1,\dots$. La TIR $y$ de cada título es la raíz de $P=\sum_k CF_k(1+y)^{-t_k}$ y su duración de Macaulay $D=\sum_k t_kCF_k(1+y)^{-t_k}/P$.

\section{Metodología}
\subsection{Modelo de Svensson}
Los tipos al contado con capitalización continua son
\begin{equation}r_T=\beta_0+(\beta_1+\beta_2)\frac{1-e^{-\kappa_1T}}{\kappa_1T}-\beta_2e^{-\kappa_1T}+\beta_3\frac{1-e^{-\kappa_2T}}{\kappa_2T}-\beta_3e^{-\kappa_2T},\end{equation}
con función de descuento $D(0,T)=e^{-r_TT}$ y forward instantáneo $f(T)=\beta_0+\beta_1e^{-\kappa_1T}+\beta_2\kappa_1Te^{-\kappa_1T}+\beta_3\kappa_2Te^{-\kappa_2T}$. $\beta_0$ es el tipo a plazo infinito, $\beta_0+\beta_1$ el tipo a plazo cero, y $\beta_2,\beta_3$ con $\kappa_1,\kappa_2$ gobiernan dos jorobas.
El precio teórico del título $j$ es $\hat P_j=\sum_kCF_{jk}\,e^{-r(t_{jk})t_{jk}}$. Siguiendo el Anexo~I, que advierte de que mínimos cuadrados ordinarios sobre precios penaliza demasiado los errores de los títulos largos (un mismo error de TIR produce un error de precio proporcional a la duración), se estiman los parámetros minimizando
\begin{equation}\min_{\beta,\kappa}\ \sum_{j=1}^{n}\Bigl(\frac{\hat P_j-P_j}{D_j}\Bigr)^2.\end{equation}
La optimización (mínimos cuadrados no lineales, algoritmo de región de confianza) parte de los valores del enunciado $(\beta_0,\beta_1,\beta_2,\beta_3,\kappa_1,\kappa_2)=(0{,}05;-0{,}025;0{,}01;-0{,}04;1{,}2;0{,}5)$ y se repite desde otros cinco arranques, con las cotas $\kappa\in[0{,}05;10]$, $|\beta_{1,2,3}|\le0{,}1$ y $\beta_0\in[0;0{,}2]$. Se queda el óptimo de menor coste.

\subsection{Splines cúbicas}
Se aproxima $D(0,T)=1+a_1T+a_2T^2+a_3T^3+\sum_{p}b_p(T-\xi_p)_+^3$ (con $D(0,0)=1$ impuesto). Es una función de clase $C^2$ y el precio de cada título es lineal en $(a_1,a_2,a_3,b_1,\dots)$, de modo que la estimación es un mínimo cuadrado ponderado (con el mismo peso $1/D_j$) resuelto con álgebra lineal. Nudos internos: $\xi=(1{,}5, 4{,}0, 8{,}0, 15{,}0)$ años, es decir 7 parámetros. Se comprueba que $D$ es positiva y decreciente en $[0;45{,}1]$ (se cumple).

\section{Resultados}
\subsection{Svensson}
\begin{table}[H]\centering\caption{Parámetros estimados (Svensson).}
\begin{tabular}{lc}\toprule Parámetro & Valor\\\midrule
$\beta_0$ & 0{,}04930 \\
$\beta_1$ & -0{,}02321 \\
$\beta_2$ & -0{,}01750 \\
$\beta_3$ & -0{,}04232 \\
$\kappa_1$ & 3{,}04691 \\
$\kappa_2$ & 0{,}46020 \\
\bottomrule\end{tabular}\end{table}
\begin{figure}[H]\centering\includegraphics[width=.95\textwidth]{fig_svensson.pdf}\caption{Curva cupón cero de Svensson (línea punteada: tipo $r_T$ con capitalización continua; línea continua: tipo anual equivalente $e^{r_T}-1$, comparable con la TIR) y TIR de los 57 bonos y obligaciones.}\label{fig:sv}\end{figure}
El ajuste (Tabla~\ref{tab:met}) deja un error medio absoluto de 2{,}3 pb en TIR; el mayor error es de 23{,}2 pb (OBL 1,450 10/2071, plazo 45{,}1 años).

\subsection{Splines cúbicas y comparación}
\begin{table}[H]\centering\caption{Calidad del ajuste (en la muestra).}\label{tab:met}
\begin{tabular}{lcc}\toprule & Svensson & Splines\\\midrule
Nº de parámetros & 6 & 7 \\
RMSE precio (€ por 100 nominal) & 0{,}380 & 0{,}254 \\
RMSE TIR (pb) & 3{,}92 & 2{,}81 \\
Error absoluto medio TIR (pb) & 2{,}31 & 2{,}05 \\
Error máximo TIR (pb) & 23{,}20 & 10{,}58 \\
Función objetivo $\sum((\hat P-P)/D)^2$ & 0{,}03082 & 0{,}02493 \\
Criterio BIC (sobre residuos ponderados) & -404{,}5 & -412{,}6 \\
\bottomrule\end{tabular}\end{table}
\begin{figure}[H]\centering\includegraphics[width=.95\textwidth]{fig_comparacion.pdf}\caption{Curvas cupón cero (tipo anual equivalente) de Svensson y de splines (líneas verticales: nudos) y TIR.}\end{figure}
\begin{table}[H]\centering\caption{Tipos al contado y forward instantáneos estimados (\%). $r$: capitalización continua; $R=e^r-1$: anual equivalente.}
\begin{tabular}{rcccccc}\toprule & \multicolumn{3}{c}{Svensson} & \multicolumn{3}{c}{Splines}\\\cmidrule(lr){2-4}\cmidrule(lr){5-7}$T$ (años) & $r_T$ & $R_T$ & $f(T)$ & $r_T$ & $R_T$ & $f(T)$\\\midrule
0{,}25 & 2{,}672 & 2{,}708 & 2{,}790 & 2{,}690 & 2{,}726 & 2{,}820 \\
0{,}50 & 2{,}803 & 2{,}843 & 3{,}069 & 2{,}811 & 2{,}851 & 3{,}038 \\
1 & 3{,}020 & 3{,}066 & 3{,}337 & 3{,}005 & 3{,}050 & 3{,}323 \\
2 & 3{,}187 & 3{,}238 & 3{,}349 & 3{,}192 & 3{,}244 & 3{,}373 \\
3 & 3{,}254 & 3{,}308 & 3{,}459 & 3{,}259 & 3{,}313 & 3{,}437 \\
5 & 3{,}431 & 3{,}491 & 3{,}954 & 3{,}425 & 3{,}485 & 3{,}959 \\
7 & 3{,}646 & 3{,}714 & 4{,}386 & 3{,}646 & 3{,}714 & 4{,}403 \\
10 & 3{,}928 & 4{,}006 & 4{,}735 & 3{,}924 & 4{,}002 & 4{,}707 \\
15 & 4{,}233 & 4{,}323 & 4{,}901 & 4{,}239 & 4{,}330 & 4{,}974 \\
20 & 4{,}404 & 4{,}502 & 4{,}926 & 4{,}428 & 4{,}527 & 4{,}998 \\
30 & 4{,}579 & 4{,}685 & 4{,}930 & 4{,}576 & 4{,}682 & 4{,}650 \\
45 & 4{,}696 & 4{,}808 & 4{,}930 & 4{,}420 & 4{,}519 & 3{,}698 \\
\bottomrule\end{tabular}\end{table}
\begin{figure}[H]\centering\includegraphics[width=.85\textwidth]{fig_forward.pdf}\caption{Forward instantáneo implícito.}\end{figure}
\begin{figure}[H]\centering\includegraphics[width=\textwidth]{fig_residuos.pdf}\caption{Error del modelo en TIR (modelo menos mercado) según el plazo del bono.}\end{figure}

\section{Comentario de los resultados}
\paragraph{Forma de la curva.} La ETTI es \emph{creciente} y cóncava: parte del 2{,}62\,\% a 0,1 años, alcanza el 3,02\,\% al año y se acerca al 4\,\% a los 10 años y se aproxima al $\beta_0=4{,}93$\,\% (cap. continua) en el tramo largo; en tipo anual equivalente la asíntota es $e^{\beta_0}-1=5{,}05$\,\%. $\beta_0+\beta_1=2{,}61$\,\% es el tipo a plazo cero. La pendiente positiva implica, bajo la Teoría Pura de las Expectativas, que se esperan subidas de tipos (forward a 1 año $\approx3{,}34$\,\%, a 5 años $\approx3{,}95$\,\%); bajo la Preferencia por la Liquidez parte de esa pendiente es prima de plazo y la subida esperada de los tipos sería menor. Es coherente con la subida del Euribor de 2026 observada en el Trabajo~1 (el 12 meses pasa del 2,73\,\% al 3,30\,\% entre julio y octubre).
\paragraph{TIR frente a tipo cupón cero.} La curva cupón cero y las TIR no coinciden porque la TIR es una media de los tipos al contado ponderada por los flujos del bono (sesgo del cupón): en una curva creciente, los bonos con cupón elevado tienen TIR inferior al cupón cero del mismo vencimiento. En la muestra los puntos se distribuyen en torno a la curva (error absoluto medio de 2{,}3 pb) con una ligera dispersión por la diferente duración de los cupones y por la liquidez de cada referencia.
\paragraph{Calidad del ajuste.} Los errores son pequeños (menos de 4 pb en media cuadrática) pero no nulos: pueden deberse a diferencias de liquidez, a asincronías entre las cotizaciones, a que se trabaja con la media de las posiciones de compra y venta y a la restricción de forma de la función elegida (el Anexo~I recuerda que los tipos de la ETTI no son observables). El mayor residuo de Svensson está en el tramo más largo (Obl.\ 1,450\,\% 10/2071, 45 años, la TIR del modelo queda por encima de la de mercado), donde hay pocos títulos, poco líquidos.
\paragraph{Svensson frente a splines.} La spline ajusta mejor en la muestra (RMSE en TIR de 2{,}8 frente a 3{,}9 pb, y también menor BIC con un parámetro más), pero su tramo largo no es fiable: el forward instantáneo baja de más del 5\,\% a 20 años hasta 3{,}7\,\% a 45 años, y el polinomio no está acotado, de modo que extrapolar más allá del último vencimiento carece de sentido (limitación señalada en el Anexo~I). Svensson impone una forma suave con forward asintótico constante ($\approx4{,}93$\,\%), más estable en el extremo largo y coherente con la teoría, a costa de un ajuste algo peor en el tramo corto y medio.
\paragraph{Robustez e identificación.} Los parámetros de Svensson están débilmente identificados: $(\beta_2,\kappa_1)$ y $(\beta_3,\kappa_2)$ describen formas parecidas y la función objetivo tiene un valle casi plano. La Tabla~\ref{tab:arr} muestra que seis arranques distintos convergen a tres combinaciones de parámetros muy diferentes con costes casi idénticos (diferencias del orden del 0,3\,\%), y la Tabla~\ref{tab:rob} que las curvas resultantes difieren menos de 4 pb en todos los plazos salvo en el de un año con mínimos cuadrados sin ponderar (unos 8 pb). Por tanto, la interpretación económica de cada $\beta$ debe hacerse con cautela, y solo la curva (no cada parámetro) es robusta.
\begin{table}[H]\centering\caption{Multiarranque de Svensson: valores iniciales y finales $(\beta_0,\beta_1,\beta_2,\beta_3,\kappa_1,\kappa_2)$ y coste $\tfrac12\sum e^2$.}\label{tab:arr}\footnotesize
\begin{tabular}{c cccccc cccccc c}\toprule \# & \multicolumn{6}{c}{inicio} & \multicolumn{6}{c}{final} & coste\\\midrule
1 & 0{,}0500 & -0{,}0250 & 0{,}0100 & -0{,}0400 & 1{,}2000 & 0{,}5000 & 0{,}0492 & -0{,}0243 & 0{,}0788 & -0{,}1000 & 0{,}7162 & 0{,}5599 & 0{,}015451 \\
2 & 0{,}0400 & 0{,}0000 & 0{,}0000 & 0{,}0000 & 1{,}0000 & 0{,}3000 & 0{,}0492 & -0{,}0243 & 0{,}0788 & -0{,}1000 & 0{,}7162 & 0{,}5599 & 0{,}015451 \\
3 & 0{,}0450 & -0{,}0200 & 0{,}0200 & -0{,}0200 & 2{,}0000 & 0{,}2000 & 0{,}0493 & -0{,}0232 & -0{,}0175 & -0{,}0423 & 3{,}0469 & 0{,}4602 & 0{,}015411 \\
4 & 0{,}0500 & -0{,}0300 & 0{,}0300 & 0{,}0300 & 0{,}5000 & 3{,}0000 & 0{,}0492 & -0{,}0243 & -0{,}1000 & 0{,}0847 & 0{,}5613 & 0{,}7121 & 0{,}015442 \\
5 & 0{,}0600 & -0{,}0400 & -0{,}0200 & 0{,}0200 & 0{,}8000 & 0{,}1000 & 0{,}0492 & -0{,}0243 & 0{,}0788 & -0{,}1000 & 0{,}7162 & 0{,}5599 & 0{,}015451 \\
6 & 0{,}0500 & -0{,}0200 & 0{,}0000 & 0{,}0000 & 2{,}5000 & 0{,}5000 & 0{,}0493 & -0{,}0232 & -0{,}0175 & -0{,}0423 & 3{,}0469 & 0{,}4602 & 0{,}015411 \\
\bottomrule\end{tabular}\end{table}
\begin{table}[H]\centering\caption{Tipo al contado $r_T$ (\%, cap. continua) con distintas especificaciones de la estimación.}\label{tab:rob}\footnotesize
\begin{tabular}{lcccc}\toprule Especificación & $T=1$ & $T=5$ & $T=10$ & $T=30$\\\midrule
Base: ponderación $1/D$, cotas $|\beta|\le0{,}1$ & 3{,}020 & 3{,}431 & 3{,}928 & 4{,}579 \\
MCO sobre precios (sin ponderar) & 3{,}099 & 3{,}425 & 3{,}946 & 4{,}542 \\
Ponderando por 1/(D\_mod·P) ($\approx$ error en TIR) & 3{,}025 & 3{,}435 & 3{,}928 & 4{,}537 \\
Sin cotas en $\beta$ (solo $\kappa$>0) & 3{,}015 & 3{,}431 & 3{,}929 & 4{,}574 \\
\bottomrule\end{tabular}\end{table}
\begin{table}[H]\centering\caption{Sensibilidad de la spline a los nudos.}\footnotesize
\begin{tabular}{lccc}\toprule Nudos (años) & Parámetros & $\sum((\hat P-P)/D)^2$ & RMSE TIR (pb)\\\midrule
2, 5, 10 & 6 & 0{,}03367 & 3{,}08 \\
1{,}5, 4, 8, 15 & 7 & 0{,}02493 & 2{,}81 \\
2, 5, 10, 20 & 7 & 0{,}03120 & 2{,}84 \\
1, 3, 6, 10, 20 & 8 & 0{,}02144 & 2{,}51 \\
2, 5, 10, 20, 30 & 8 & 0{,}03067 & 2{,}79 \\
\bottomrule\end{tabular}\end{table}

\section{Conclusiones}
\begin{enumerate}
\item La ETTI cupón cero del Tesoro el 11/09/2026 es creciente y cóncava, del 2{,}6\,\% a plazos cortos al 4{,}7\,\% a 45 años (cap. continua).
\item El modelo de Svensson reproduce las TIR de los 57 títulos con un error cuadrático medio de 3{,}9 pb, con un tramo largo estable; las splines ajustan algo mejor (2{,}8 pb) pero extrapolan mal.
\item Los parámetros individuales de Svensson no están bien identificados (valle plano), aunque la curva sí lo está.
\item La diferencia entre TIR y tipos cupón cero refleja el sesgo del cupón y la distinta liquidez de los títulos.
\end{enumerate}

\appendix
\section{Datos y ajuste por título}
\footnotesize
\begin{longtable}{lcccccccccc}
\toprule Título & Cupón & Vencim. & $P^*$ & $P$ bruto & $T$ & $D$ & TIR (\%) & TIR Sv. & TIR Spl.\\\midrule\endhead
BON 0,000 01/2027 & 0{,}00 & 31/01/2027 & 98{,}972 & 98{,}972 & 0{,}38 & 0{,}38 & 2{,}749 & 2{,}777 & 2{,}794 \\
BON 2,350 03/2029 & 2{,}35 & 30/03/2029 & 97{,}683 & 98{,}764 & 2{,}54 & 2{,}47 & 3{,}312 & 3{,}273 & 3{,}281 \\
BON 2,400 05/2028 & 2{,}40 & 31/05/2028 & 98{,}672 & 99{,}369 & 1{,}71 & 1{,}69 & 3{,}206 & 3{,}207 & 3{,}209 \\
BON 2,500 05/2027 & 2{,}50 & 31/05/2027 & 99{,}653 & 100{,}379 & 0{,}71 & 0{,}71 & 2{,}991 & 2{,}950 & 2{,}943 \\
BON 2,600 05/2031 & 2{,}60 & 31/05/2031 & 96{,}240 & 96{,}995 & 4{,}71 & 4{,}46 & 3{,}477 & 3{,}451 & 3{,}445 \\
BON 2,700 01/2030 & 2{,}70 & 31/01/2030 & 97{,}915 & 99{,}587 & 3{,}38 & 3{,}22 & 3{,}359 & 3{,}331 & 3{,}332 \\
BON 3,500 05/2029 & 3{,}50 & 31/05/2029 & 100{,}507 & 101{,}524 & 2{,}71 & 2{,}61 & 3{,}297 & 3{,}283 & 3{,}290 \\
OBL 0,000 01/2028 & 0{,}00 & 31/01/2028 & 95{,}833 & 95{,}833 & 1{,}38 & 1{,}38 & 3{,}131 & 3{,}161 & 3{,}155 \\
OBL 0,100 04/2031 & 0{,}10 & 30/04/2031 & 85{,}928 & 85{,}965 & 4{,}62 & 4{,}61 & 3{,}444 & 3{,}451 & 3{,}444 \\
OBL 0,500 04/2030 & 0{,}50 & 30/04/2030 & 90{,}438 & 90{,}625 & 3{,}62 & 3{,}59 & 3{,}345 & 3{,}355 & 3{,}353 \\
OBL 0,500 10/2031 & 0{,}50 & 31/10/2031 & 86{,}213 & 86{,}648 & 5{,}13 & 5{,}05 & 3{,}482 & 3{,}503 & 3{,}497 \\
OBL 0,600 10/2029 & 0{,}60 & 31/10/2029 & 92{,}075 & 92{,}598 & 3{,}13 & 3{,}09 & 3{,}308 & 3{,}316 & 3{,}320 \\
OBL 0,700 04/2032 & 0{,}70 & 30/04/2032 & 85{,}750 & 86{,}013 & 5{,}62 & 5{,}51 & 3{,}538 & 3{,}556 & 3{,}552 \\
OBL 0,800 07/2027 & 0{,}80 & 30/07/2027 & 98{,}115 & 98{,}216 & 0{,}87 & 0{,}87 & 3{,}016 & 3{,}021 & 3{,}007 \\
OBL 0,800 07/2029 & 0{,}80 & 30/07/2029 & 93{,}195 & 93{,}296 & 2{,}87 & 2{,}85 & 3{,}321 & 3{,}298 & 3{,}304 \\
OBL 0,850 07/2037 & 0{,}85 & 30/07/2037 & 72{,}347 & 72{,}455 & 10{,}87 & 10{,}32 & 4{,}043 & 4{,}057 & 4{,}052 \\
OBL 1,000 07/2042 & 1{,}00 & 30/07/2042 & 62{,}710 & 62{,}836 & 15{,}87 & 14{,}34 & 4{,}286 & 4{,}324 & 4{,}333 \\
OBL 1,000 10/2050 & 1{,}00 & 31/10/2050 & 49{,}550 & 50{,}421 & 24{,}13 & 19{,}76 & 4{,}454 & 4{,}534 & 4{,}560 \\
OBL 1,200 10/2040 & 1{,}20 & 31/10/2040 & 68{,}495 & 69{,}540 & 14{,}13 & 12{,}61 & 4{,}202 & 4{,}242 & 4{,}244 \\
OBL 1,250 10/2030 & 1{,}25 & 31/10/2030 & 91{,}893 & 92{,}982 & 4{,}13 & 4{,}00 & 3{,}387 & 3{,}398 & 3{,}392 \\
OBL 1,300 10/2026 & 1{,}30 & 31/10/2026 & 99{,}825 & 100{,}958 & 0{,}13 & 0{,}13 & 2{,}664 & 2{,}659 & 2{,}660 \\
OBL 1,400 04/2028 & 1{,}40 & 30/04/2028 & 97{,}190 & 97{,}715 & 1{,}62 & 1{,}61 & 3{,}199 & 3{,}198 & 3{,}198 \\
OBL 1,400 07/2028 & 1{,}40 & 30/07/2028 & 96{,}727 & 96{,}904 & 1{,}87 & 1{,}86 & 3{,}226 & 3{,}225 & 3{,}230 \\
OBL 1,450 04/2029 & 1{,}45 & 30/04/2029 & 95{,}505 & 96{,}049 & 2{,}62 & 2{,}58 & 3{,}263 & 3{,}280 & 3{,}287 \\
OBL 1,450 10/2027 & 1{,}45 & 31/10/2027 & 98{,}210 & 99{,}473 & 1{,}13 & 1{,}11 & 3{,}086 & 3{,}103 & 3{,}090 \\
OBL 1,450 10/2071 & 1{,}45 & 31/10/2071 & 42{,}305 & 43{,}568 & 45{,}13 & 25{,}12 & 4{,}424 & 4{,}656 & 4{,}478 \\
OBL 1,500 04/2027 & 1{,}50 & 30/04/2027 & 99{,}130 & 99{,}693 & 0{,}62 & 0{,}62 & 2{,}917 & 2{,}909 & 2{,}907 \\
OBL 1,850 07/2035 & 1{,}85 & 30/07/2035 & 85{,}285 & 85{,}518 & 8{,}87 & 8{,}18 & 3{,}837 & 3{,}879 & 3{,}878 \\
OBL 1,900 10/2052 & 1{,}90 & 31/10/2052 & 60{,}745 & 62{,}400 & 26{,}13 & 18{,}49 & 4{,}479 & 4{,}521 & 4{,}541 \\
OBL 2,350 07/2033 & 2{,}35 & 30/07/2033 & 92{,}150 & 92{,}446 & 6{,}87 & 6{,}39 & 3{,}662 & 3{,}680 & 3{,}679 \\
OBL 2,550 10/2032 & 2{,}55 & 31/10/2032 & 94{,}395 & 96{,}617 & 6{,}13 & 5{,}61 & 3{,}584 & 3{,}600 & 3{,}597 \\
OBL 2,700 10/2048 & 2{,}70 & 31/10/2048 & 76{,}120 & 78{,}472 & 22{,}13 & 15{,}57 & 4{,}412 & 4{,}433 & 4{,}455 \\
OBL 2,900 10/2046 & 2{,}90 & 31/10/2046 & 80{,}595 & 83{,}122 & 20{,}13 & 14{,}46 & 4{,}368 & 4{,}388 & 4{,}407 \\
OBL 3,000 01/2033 & 3{,}00 & 31/01/2033 & 96{,}430 & 98{,}288 & 6{,}38 & 5{,}78 & 3{,}635 & 3{,}624 & 3{,}622 \\
OBL 3,100 07/2031 & 3{,}10 & 30/07/2031 & 98{,}328 & 98{,}718 & 4{,}87 & 4{,}58 & 3{,}478 & 3{,}466 & 3{,}460 \\
OBL 3,150 04/2033 & 3{,}15 & 30/04/2033 & 97{,}115 & 98{,}297 & 6{,}62 & 6{,}00 & 3{,}646 & 3{,}648 & 3{,}647 \\
OBL 3,150 04/2035 & 3{,}15 & 30/04/2035 & 95{,}047 & 96{,}230 & 8{,}62 & 7{,}57 & 3{,}833 & 3{,}841 & 3{,}840 \\
OBL 3,200 10/2035 & 3{,}20 & 31/10/2035 & 94{,}885 & 97{,}673 & 9{,}13 & 7{,}80 & 3{,}875 & 3{,}883 & 3{,}882 \\
OBL 3,250 04/2034 & 3{,}25 & 30/04/2034 & 96{,}722 & 97{,}942 & 7{,}62 & 6{,}78 & 3{,}750 & 3{,}747 & 3{,}748 \\
OBL 3,300 04/2036 & 3{,}30 & 30/04/2036 & 95{,}118 & 96{,}356 & 9{,}62 & 8{,}26 & 3{,}917 & 3{,}922 & 3{,}919 \\
OBL 3,400 10/2036 & 3{,}40 & 31/10/2036 & 95{,}412 & 98{,}375 & 10{,}13 & 8{,}44 & 3{,}958 & 3{,}959 & 3{,}955 \\
OBL 3,450 07/2043 & 3{,}45 & 30/07/2043 & 89{,}713 & 90{,}147 & 16{,}87 & 12{,}76 & 4{,}321 & 4{,}289 & 4{,}299 \\
OBL 3,450 07/2066 & 3{,}45 & 30/07/2066 & 80{,}295 & 80{,}730 & 39{,}87 & 20{,}09 & 4{,}526 & 4{,}566 & 4{,}492 \\
OBL 3,450 10/2034 & 3{,}45 & 31/10/2034 & 97{,}670 & 100{,}676 & 8{,}13 & 7{,}00 & 3{,}787 & 3{,}793 & 3{,}794 \\
OBL 3,500 01/2041 & 3{,}50 & 31/01/2041 & 92{,}188 & 94{,}355 & 14{,}38 & 11{,}15 & 4{,}235 & 4{,}197 & 4{,}199 \\
OBL 3,550 10/2033 & 3{,}55 & 31/10/2033 & 99{,}083 & 102{,}175 & 7{,}13 & 6{,}23 & 3{,}698 & 3{,}697 & 3{,}697 \\
OBL 3,900 07/2039 & 3{,}90 & 30/07/2039 & 97{,}620 & 98{,}112 & 12{,}87 & 10{,}28 & 4{,}141 & 4{,}120 & 4{,}117 \\
OBL 3,950 10/2056 & 3{,}95 & 31/10/2056 & 90{,}500 & 93{,}941 & 30{,}13 & 16{,}85 & 4{,}534 & 4{,}497 & 4{,}501 \\
OBL 4,000 10/2054 & 4{,}00 & 31/10/2054 & 92{,}105 & 95{,}590 & 28{,}13 & 16{,}29 & 4{,}500 & 4{,}478 & 4{,}490 \\
OBL 4,200 01/2037 & 4{,}20 & 31/01/2037 & 101{,}843 & 104{,}443 & 10{,}38 & 8{,}43 & 3{,}978 & 3{,}965 & 3{,}961 \\
OBL 4,250 07/2047 & 4{,}25 & 30/07/2047 & 98{,}165 & 98{,}701 & 20{,}87 & 14{,}10 & 4{,}385 & 4{,}370 & 4{,}388 \\
OBL 4,700 07/2041 & 4{,}70 & 30/07/2041 & 105{,}227 & 105{,}820 & 14{,}87 & 11{,}08 & 4{,}219 & 4{,}194 & 4{,}198 \\
OBL 4,900 07/2040 & 4{,}90 & 30/07/2040 & 107{,}567 & 108{,}185 & 13{,}87 & 10{,}47 & 4{,}170 & 4{,}150 & 4{,}150 \\
OBL 5,150 10/2028 & 5{,}15 & 31/10/2028 & 103{,}868 & 108{,}354 & 2{,}13 & 1{,}99 & 3{,}236 & 3{,}243 & 3{,}249 \\
OBL 5,150 10/2044 & 5{,}15 & 31/10/2044 & 110{,}280 & 114{,}767 & 18{,}13 & 12{,}02 & 4{,}320 & 4{,}289 & 4{,}302 \\
OBL 5,750 07/2032 & 5{,}75 & 30/07/2032 & 111{,}560 & 112{,}285 & 5{,}87 & 5{,}16 & 3{,}534 & 3{,}558 & 3{,}554 \\
OBL 6,000 01/2029 & 6{,}00 & 31/01/2029 & 106{,}185 & 109{,}900 & 2{,}38 & 2{,}22 & 3{,}248 & 3{,}258 & 3{,}265 \\
\bottomrule\end{longtable}
\normalsize
\section{Archivos y código}
El libro \texttt{Trabajo2\_ETTI\_RentaFija.xlsx} reproduce todo el cálculo con fórmulas (hojas \emph{Parametros}, \emph{Bonos}, \emph{Flujos}, \emph{Splines}, \emph{Curva}); en \emph{Parametros} está preparada la función objetivo de Svensson para Solver y los coeficientes de las splines se obtienen con \texttt{LINEST}.
\subsection*{t2core.py}
\lstinputlisting[language=Python]{listing_t2core.py}
\subsection*{t2.py}
\lstinputlisting[language=Python]{listing_t2.py}
\end{document}
